\subsection{Differential Forms and Manifolds}

Differential forms package oriented integration on parametrized cells. They use determinant sign and laws (Definition \ref{def:matrix-trace-and-determinant}, Theorems \ref{thm:determinant-basic-laws}, \ref{thm:cauchy-binet-minors}) and change of variables (Theorem \ref{thm:change-of-variables-rk}); conceptually, forms test oriented cells.

\begin{perspective} \label{per:differential-forms-as-functionals} \hlt{Forms as Functionals}\\
The notation $\int f\,dx$ treats $dx$ as directed coordinate variation. Differential forms make this precise: a $1$-form such as $f\,dx+g\,dy$ assigns a number to an oriented path; $dx$ records net $x$-variation, and $f\,dx$ weights it by $f$.
\end{perspective}

\begin{figure}[H]
  \centering
  \begin{tikzpicture}[
  scale=0.9,
  >=Stealth,
  every node/.style={font=\scriptsize},
  axis/.style={black!40, line width=0.5pt},
  positive/.style={
    blue!70!black,
    line width=0.85pt,
    postaction={
      decorate,
      decoration={
        markings,
        mark=at position 0.18 with {\arrow{Stealth[length=4pt,width=4pt]}},
        mark=at position 0.45 with {\arrow{Stealth[length=4pt,width=4pt]}},
        mark=at position 0.76 with {\arrow{Stealth[length=4pt,width=4pt]}}
      }
    }
  },
  negative/.style={
    red!70!black,
    line width=0.85pt,
    postaction={
      decorate,
      decoration={
        markings,
        mark=at position 0.22 with {\arrow{Stealth[length=4pt,width=4pt]}},
        mark=at position 0.55 with {\arrow{Stealth[length=4pt,width=4pt]}},
        mark=at position 0.82 with {\arrow{Stealth[length=4pt,width=4pt]}}
      }
    }
  },
  point/.style={circle, fill=black, inner sep=1.15pt}
]
  \draw[axis] (-0.35,0) -- (6.45,0) node[right] {$x$};
  \draw[axis] (0,-1.55) -- (0,1.65) node[above] {$y$};

  \node[point] (p) at (0.45,0) {};
  \node[point] (q) at (5.85,0) {};
  \node[below left=2pt] at (p) {$p$};
  \node[below right=2pt] at (q) {$q$};

  \draw[positive]
    (0.45,0)
    .. controls (1.20,0.70) and (2.10,0.45) .. (2.80,0.70)
    .. controls (3.35,0.95) and (3.70,0.05) .. (4.20,0.28)
    .. controls (4.85,0.58) and (5.35,0.34) .. (5.85,0);
  \node[blue!70!black] at (4.55,0.86) {$C,\ y>0$};

  \draw[negative]
    (0.45,0)
    .. controls (1.15,-0.72) and (2.35,-1.03) .. (3.30,-0.86)
    .. controls (4.35,-0.66) and (5.10,-0.28) .. (5.85,0);
  \node[red!70!black] at (3.55,-1.18) {$C',\ y<0$};
\end{tikzpicture}

  \caption{$y\,dx$ records signed, $y$-weighted $x$-variation.}
  \label{fig:forms-weighted-variation-paths}
\end{figure}

\begin{definition} \label{def:forms-manifolds-core-source-todo}\label{def:k-surfaces-and-differential-forms} \hlt{$k$-Surfaces and Differential Forms}\\
Let $E\subseteq\R^n$ be open.
\begin{enumerate}[label=(\roman*)]
\setlength{\itemsep}{0pt}
\item A compact \hlt{parameter domain} $D\subseteq\R^k$ is a $k$-cell or $Q^k=\{\mathbf{u}\in\R^k:u_i\geq0,\ \sum\limits_i u_i\leq1\}$.\\
For compact $D$, $C^r(D)$ means extendable as $C^r$ to an open neighborhood of $D$.
\item A \hlt{$k$-surface} in $E$ is a $C^1$ map $\Phi:D\to E$ from a compact parameter domain $D\subseteq\R^k$.
\item For $k\geq1$, a \hlt{$k$-form} on $E$ is $\omega=\sum_I a_I(\x)\,dx_{i_1}\wedge\cdots\wedge dx_{i_k}$, where $I=(i_1,\ldots,i_k)$ and $a_I\in C(E,\R)$.
\item For an increasing $k$-index $I=(i_1<\cdots<i_k)$, set $dx_I\defeq dx_{i_1}\wedge\cdots\wedge dx_{i_k}$.\\
A \hlt{standard presentation} is $\omega=\sum_I b_I(\x)\,dx_I$, with $I$ increasing.
\item For a $k$-surface $\Phi:D\to E$, define $\int_{\Phi}\omega\defeq\int_D\sum_I a_I(\Phi(\mathbf{u}))J_I\Phi(\mathbf{u})\,d\mathbf{u}$, where $J_I\Phi\defeq\frac{\partial(x_{i_1},\ldots,x_{i_k})}{\partial(u_1,\ldots,u_k)}$.
\item A $0$-form is a continuous function on $E$; a form is $C^r$ iff all coefficients $a_I$ are $C^r$.
\end{enumerate}
\end{definition}

\begin{remark} \label{rem:parametrized-cells-not-images} \hlt{Parametrized Cells, Not Just Images}\\
Parametrizations may fail to be injective or regular. Forms therefore act on parametrized cells, not just images; this matters for orientation, multiplicity, and cancellation.
\end{remark}

\begin{proposition} \label{prop:basic-forms-standard-presentation} \hlt{Basic Forms and Standard Presentation}
\begin{enumerate}[label=(\roman*)]
\setlength{\itemsep}{0pt}
\item Swapping neighboring factors in $dx_{i_1}\wedge\cdots\wedge dx_{i_k}$ changes sign; repeated indices give $0$.
\item Every $k$-form has a unique standard presentation.
\end{enumerate}
\begin{proof}
\begin{enumerate}[label=(\roman*)]
\setlength{\itemsep}{0pt}
\item For $\Phi(\mathbf{u})=(x_{i_1},\ldots,x_{i_k})$, the coefficient of $dx_{i_1}\wedge\cdots\wedge dx_{i_k}$ is the determinant of the corresponding $k\times k$ Jacobian minor. Row swaps change its sign and repeated rows make it $0$.
\item By (i), each term either vanishes or sorts to an increasing index, giving existence.\\
For uniqueness, let $\omega=\sum_I b_I(\x)\,dx_I$ be zero in standard presentation.\\
If $b_J(\mathbf{v})\neq0$ for increasing $J=(j_1,\ldots,j_k)$, choose $D=[-h,h]^k$ so $b_J(\Phi(\mathbf{u}))$ has constant nonzero sign for $\Phi(\mathbf{u})=\mathbf{v}+\sum\limits_{r=1}^{k}u_r\mathbf{e}_{j_r}$. Then the $J$-Jacobian is $1$ and all other standard Jacobians are $0$, so $0=\int_{\Phi}\omega=\int_D b_J(\Phi(\mathbf{u}))\,d\mathbf{u}$, a contradiction. Hence all $b_I=0$.
\end{enumerate}
\hlt{Structural Commentary}: Wedge signs are determinant signs.\\
\hlt{Key Insight}: Standard form sorts indices once.
\end{proof}
\end{proposition}

\begin{theorem} \label{thm:forms-reparametrization-naturality} \hlt{Naturality Under Reparametrization}\\
Let $\Phi:D\to E\subseteq\R^n$ be a $C^1$ $k$-surface, and let $R:D'\to D$ be a $C^1$ diffeomorphism with $\epsilon\defeq\operatorname{sgn}\det DR$ constant on $D'$. For every $k$-form $\omega$ on $E$, $\int_{\Phi\circ R}\omega=\epsilon\int_{\Phi}\omega$.
\begin{proof}
By linearity, take $\omega=f\,dx_I$. Cauchy-Binet gives $J_I(\Phi\circ R)(\mathbf{u})=J_I\Phi(R(\mathbf{u}))\det DR(\mathbf{u})$.\\
Thus $\int_{\Phi\circ R}\omega=\int_{D'}f(\Phi(R\mathbf{u}))J_I\Phi(R\mathbf{u})\det DR(\mathbf{u})\,d\mathbf{u}=\epsilon\int_D f(\Phi(\mathbf{v}))J_I\Phi(\mathbf{v})\,d\mathbf{v}$ by change of variables.\\
\hlt{Structural Commentary}: Reparametrization changes coordinates only.\\
\hlt{Key Insight}: Orientation contributes $\operatorname{sgn}\det DR$.
\end{proof}
\end{theorem}

\begin{definition} \label{def:wedge-product-and-exterior-derivative} \hlt{Product and Exterior Derivative}\\
For basic forms, $dx_I\wedge dx_J$ is obtained by concatenating $I,J$; it is $0$ if an index repeats, and otherwise $\operatorname{sgn}(\sigma)dx_K$, where $K$ is the increasing rearrangement and $\sigma$ is the sorting permutation. Extend bilinearly.
\begin{enumerate}[label=(\roman*)]
\setlength{\itemsep}{0pt}
\item For $f\in C^1(E)$, define $df\defeq\sum\limits_{i=1}^{n}(D_if)(\x)\,dx_i$.
\item For a $C^1$ $k$-form $\omega=\sum_I b_I(\x)\,dx_I$, define $d\omega\defeq\sum\limits_I (db_I)\wedge dx_I$.
\end{enumerate}
\end{definition}

\begin{theorem} \label{thm:exterior-derivative-basic-laws} \hlt{Exterior Derivative Laws}\\
If $\omega,\lambda$ are $C^1$ forms in $E$ and $\omega$ is a $k$-form, then $d(\omega\wedge\lambda)=(d\omega)\wedge\lambda+(-1)^k\omega\wedge d\lambda$. If $\omega$ is $C^2$, then $d^2\omega=0$.
\begin{proof}
By bilinearity, it suffices to take $\omega=f\,dx_I$ and $\lambda=g\,dx_J$.\\
The ordinary product rule gives $d(fg)=f\,dg+g\,df$; moving $dg$ past $k$ factors in $dx_I$ gives the displayed sign.\\
For a $C^2$ $0$-form $f$, $d^2f=\sum\limits_{i,j}(D_{ij}f)(\x)\,dx_i\wedge dx_j=0$, by Theorem \ref{thm:equality-of-mixed-partials} and antisymmetry.\\
The general case follows from $\omega=\sum_I b_I\,dx_I$ and the product rule just proved.\\
\hlt{Structural Commentary}: The exterior derivative is built to make mixed partial cancellation visible.\\
\hlt{Key Insight}: Symmetric second derivatives meet antisymmetric wedge products, so $d^2=0$.
\end{proof}
\end{theorem}

\begin{definition} \label{def:pullback-of-differential-form} \hlt{Change of Variables for Forms}\\
Let $E\subseteq\R^n$, $V\subseteq\R^m$, let $T=(t_1,\ldots,t_m):E\to V$ be $C^1$, and let $\omega=\sum_I b_I(\y)\,dy_I$ be a $k$-form on $V$.
\begin{enumerate}[label=(\roman*)]
\setlength{\itemsep}{0pt}
\item Set $dt_i\defeq\sum\limits_{j=1}^{n}(D_jt_i)(\x)\,dx_j$ for $1\leq i\leq m$.
\item The \hlt{pullback} of $\omega$ by $T$ is $T^\ast\omega\defeq\omega_T\defeq\sum\limits_I b_I(T(\x))\,dt_{i_1}\wedge\cdots\wedge dt_{i_k}$.
\end{enumerate}
\end{definition}

Pullback is the differential-form analogue of the dual map in Definition \ref{def:dual-map}: cells move forward by composition with $T$, while forms move backward so they can be evaluated before or after that forward motion.

\begin{figure}[H]
  \centering
  \begin{tikzpicture}[
  scale=0.92,
  >=Stealth,
  every node/.style={font=\scriptsize},
  box/.style={draw=black!45, fill=blue!5, rounded corners=2pt, minimum width=2.25cm, minimum height=0.62cm},
  arrow/.style={->, line width=0.75pt},
  dual/.style={<-, line width=0.75pt, blue!70!black},
  pairing/.style={draw=black!40, fill=gray!6, rounded corners=2pt, inner sep=4pt}
]
  \node[box] (celln) at (0,1.15) {$k$-cell $\varphi$ in $\R^n$};
  \node[box] (cellm) at (5.2,1.15) {$T_\ast\varphi=T\circ\varphi$};
  \draw[arrow] (celln) -- node[above=2pt] {$T_\ast$} (cellm);

  \node[box] (formn) at (0,-1.15) {$T^\ast\alpha$ on $\R^n$};
  \node[box] (formm) at (5.2,-1.15) {$k$-form $\alpha$ on $\R^m$};
  \draw[dual] (formn) -- node[above=2pt] {$T^\ast$} (formm);

  \node[pairing] at (2.6,0)
    {$\displaystyle \int_{\varphi} T^\ast\alpha
      = \int_{T_\ast\varphi}\alpha$};
\end{tikzpicture}

  \caption{Cells push forward with $T$; forms pull back with $T^\ast$, so $T^\ast\alpha(\varphi)=\alpha(T_\ast\varphi)$.}
  \label{fig:forms-pushforward-pullback-duality}
\end{figure}

\begin{theorem} \label{thm:pullback-basic-laws} \hlt{Pullback Commutes with the Basic Operations}\\
Let $T=(t_1,\ldots,t_m):E\subseteq\R^n\to V\subseteq\R^m$ be $C^1$, and let $\omega_T\defeq T^\ast\omega$.
\begin{enumerate}[label=(\roman*)]
\setlength{\itemsep}{0pt}
\item $(\omega+\lambda)_T=\omega_T+\lambda_T$ for forms of the same order;
\item $(\omega\wedge\lambda)_T=\omega_T\wedge\lambda_T$;
\item $d(\omega_T)=(d\omega)_T$ when $\omega$ is $C^1$ and $T$ is $C^2$;
\item if $S:V\to W$ is $C^1$, then $(\omega_S)_T=\omega_{S\circ T}$.
\end{enumerate}
Moreover, if $\Phi$ is a $k$-surface in $E$ and $\omega$ is a $k$-form in $V$, then $\int_{T\circ\Phi}\omega=\int_{\Phi}\omega_T$.
\begin{proof}
\begin{enumerate}[label=(\roman*)]
\setlength{\itemsep}{0pt}
\item Addition follows termwise from $b_I\mapsto b_I\circ T$ and $dy_i\mapsto dt_i$.
\item Pullback replaces each $dy_i$ by $dt_i$; bilinearity gives $(\omega\wedge\lambda)_T=\omega_T\wedge\lambda_T$.
\item Let $\omega=\sum_I b_I\,dy_I$. Since $d(b_I\circ T)=(db_I)_T$ (Theorem \ref{thm:chain-rule-several-variables}) and $d(dt_i)=d^2t_i=0$, $d(\omega_T)=(d\omega)_T$.
\item If $S=(s_1,\ldots,s_p)$, both $(\omega_S)_T$ and $\omega_{S\circ T}$ replace $dz_q$ by $d(s_q\circ T)$, by the chain rule.
\end{enumerate}
For the integral identity, expand over $D$ and use Cauchy-Binet on $D(T\circ\Phi)=DT(\Phi)D\Phi$.\\
\hlt{Structural Commentary}: Pullback translates forms before integration.\\
\hlt{Key Insight}: Forms are defined so that composition of parameterizations becomes an algebraic substitution.
\end{proof}
\end{theorem}
